\chapter{Literature Review}

This chapter surveys the background and related work necessary to contextualise GraphHARE.
We cover graph neural network architectures, the role of homophily in graph learning,
reinforcement learning applied to graph problems, graph topology optimisation methods, and
position GraphHARE against its closest prior work, GraphRARE.

% -------------------------------------------------------
\section{Graph Neural Networks}
% -------------------------------------------------------

Graph Neural Networks (GNNs) generalise deep learning to graph-structured data by learning
node representations through iterative aggregation of neighbourhood information. The core
message-passing paradigm \cite{hamilton2017graphsage} updates each node's embedding by
aggregating features from its neighbours, allowing the network to capture local structural
patterns.

\subsection{Spectral Methods: GCN}

The Graph Convolutional Network (GCN) of Kipf and Welling~\cite{kipf2017gcn} defines graph
convolution in the spectral domain via a first-order approximation of spectral graph filters.
Given the symmetrically normalised adjacency $\tilde{A} = \hat{D}^{-1/2}\hat{A}\hat{D}^{-1/2}$
where $\hat{A} = A + I$, a GCN layer computes:
\begin{equation}
    H^{(l+1)} = \sigma\!\left(\tilde{A} H^{(l)} W^{(l)}\right),
\end{equation}
where $H^{(l)}$ is the node feature matrix at layer $l$ and $W^{(l)}$ is a learnable
weight matrix. GCN is transductive (requires the full graph at training time) and assumes
homophily---nodes aggregate from all neighbours equally regardless of class alignment.

\subsection{Attention-Based Methods: GAT}

Graph Attention Networks (GAT) \cite{velickovic2018gat} address the uniform aggregation
limitation of GCN by learning edge-wise attention coefficients:
\begin{equation}
    \alpha_{ij} = \frac{\exp\!\left(\text{LeakyReLU}\left(\mathbf{a}^\top
    [W\mathbf{h}_i \| W\mathbf{h}_j]\right)\right)}
    {\sum_{k \in \mathcal{N}(i)} \exp\!\left(\text{LeakyReLU}\left(\mathbf{a}^\top
    [W\mathbf{h}_i \| W\mathbf{h}_k]\right)\right)}.
\end{equation}
Neighbours that are more informative receive higher weight. GAT is still transductive and,
like GCN, can suffer in heterophilic settings where high-weight neighbours may belong to
different classes.

\subsection{Inductive Methods: GraphSAGE}

GraphSAGE (Graph Sample and Aggregate) \cite{hamilton2017graphsage} enables inductive
learning by sampling a fixed-size neighbourhood and aggregating with a learned aggregator
function (mean, LSTM, or pooling). This allows generalisation to unseen nodes and scales
to large graphs. GraphHARE uses GraphSAGE as its frozen backbone: the classifier is trained
once on the original graph and held fixed while the RL agent edits the topology.

\subsection{Advanced Propagation: APPNP}

APPNP \cite{gasteiger2019appnp} decouples feature transformation from propagation by
applying a personalised PageRank diffusion after a standard MLP:
\begin{equation}
    Z = \alpha \left(I - (1-\alpha)\tilde{A}\right)^{-1} H,
\end{equation}
where $\alpha$ is the teleportation probability. This multi-scale propagation allows APPNP
to leverage long-range neighbourhood information while retaining locality through the
teleportation term, making it competitive on both homophilic and mildly heterophilic graphs.

\subsection{Deep Residual GNNs: GCNII}

GCNII \cite{chen2020gcnii} overcomes over-smoothing in deep GCNs via two modifications:
initial residual connections (retaining a fraction of the input feature at every layer) and
identity mapping (scaling the weight matrix toward the identity). Together these allow
GCNII to stack up to 64 layers without performance collapse, capturing very long-range
dependencies.

\subsection{Transformer-Based Methods: GraphTransformer}

GraphTransformer \cite{shi2021graphtrans} replaces local neighbourhood aggregation with
global self-attention over all node pairs, treating the graph as a fully-connected
structure modulated by edge features. This enables the model to capture non-local
interactions but at quadratic computational cost in the number of nodes.

\subsection{Non-Graph Baseline: MLP}

A multilayer perceptron (MLP) trained on node features alone, ignoring graph structure,
serves as the lower-bound baseline. The gap between MLP and GNN performance quantifies
the utility of the graph topology for a given dataset. In this thesis, the MLP--GNN gap is
used as a diagnostic: it shows whether the original graph helps classification before any
editing is attempted.

% -------------------------------------------------------
\section{Homophily in Graphs}
% -------------------------------------------------------

\textit{Homophily} refers to the tendency of nodes to form edges with structurally or
semantically similar nodes. In the context of node classification, homophily is typically
measured by the \textit{node homophily ratio}:
\begin{equation}
    h = \frac{1}{|V|} \sum_{v \in V}
    \frac{\left|\{u \in \mathcal{N}(v) : y_u = y_v\}\right|}{|\mathcal{N}(v)|},
    \label{eq:homophily}
\end{equation}
where $y_v$ denotes the class label of node $v$ and $\mathcal{N}(v)$ its neighbours.
A value of $h = 1$ indicates perfect homophily (all neighbours share the same class);
$h = 0$ indicates perfect heterophily.

\subsection{Homophily and GNN Performance}

Zhu et al.~\cite{zhu2020homophily} provide a systematic analysis showing that standard
GNNs (GCN, GAT) degrade significantly on heterophilic benchmarks, while methods that
either limit aggregation depth or incorporate non-local signals are more robust. The
intuition is straightforward: in homophilic graphs, aggregating neighbour features
reinforces class-discriminative information; in heterophilic graphs, it introduces noise.
GraphHARE tests this relationship directly by treating $\Delta h$ as part of the reward
signal. This makes homophily an explicit optimisation target, while the experiments in
Chapter~4 test whether that target produces test-transferable graph improvements.

\subsection{Homophily as a Reward Signal}

To our knowledge, prior RL-based graph editing methods have not treated homophily as an
explicit reward component. GraphRARE \cite{graphrare2024} uses a relative entropy measure
to select informative edges but does not decompose the reward into homophily and task
components. GraphHARE makes homophily an explicit, tuneable objective in RL-based graph
editing, allowing the practitioner to control the structural vs.\ task-performance
trade-off via $\beta$.

% -------------------------------------------------------
\section{Graph Autoencoders}
% -------------------------------------------------------

Variational Graph Autoencoders (VGAE) \cite{kipf2016vgae} learn latent node embeddings
by encoding the graph with a GCN and decoding via inner product:
\begin{equation}
    \hat{A} = \sigma(ZZ^\top), \quad Z = \text{GCN}(X, A).
\end{equation}
The model is trained to reconstruct the adjacency matrix under a KL regularisation on
the latent space. VGAE optimises an unsupervised reconstruction objective and does not
directly target downstream task performance. In contrast, GraphHARE's reward is designed
to combine task feedback ($\Delta F_1$) with structural quality ($\Delta h$), making it a
supervised graph-editing objective rather than a reconstruction objective.

% -------------------------------------------------------
\section{Reinforcement Learning for Graph Problems}
% -------------------------------------------------------

Reinforcement learning has been applied to a range of combinatorial graph problems,
including molecular generation \cite{you2018gcpn}, network design, and graph structure
learning. The general formulation casts the graph editing process as a Markov Decision
Process (MDP): the state encodes the current graph, the action space consists of edge
additions or removals, and the reward reflects task quality or structural properties.

\subsection{Proximal Policy Optimisation}

PPO \cite{schulman2017ppo} is an on-policy actor-critic algorithm that stabilises training
by clipping the policy gradient update:
\begin{equation}
    \mathcal{L}^{\text{CLIP}}(\theta) = \mathbb{E}_t \left[
    \min\!\left(r_t(\theta)\hat{A}_t,\;
    \text{clip}(r_t(\theta), 1-\epsilon, 1+\epsilon)\hat{A}_t\right)\right],
\end{equation}
where $r_t(\theta) = \pi_\theta(a_t|s_t)/\pi_{\theta_\text{old}}(a_t|s_t)$ is the
probability ratio and $\hat{A}_t$ is the generalised advantage estimate. PPO is well-suited
to discrete action spaces with large but structured action sets, making it a natural
choice for edge-level graph editing. GraphHARE uses PPO with a GNN-based policy network
that embeds the current graph state via GraphSAGE message passing.

% -------------------------------------------------------
\section{Graph Topology Optimisation}
% -------------------------------------------------------

Graph topology optimisation encompasses methods that modify edge structure to improve
downstream task performance. Approaches include differentiable graph structure learning
(jointly optimise adjacency and GNN parameters), adversarial perturbation defences
(remove adversarial edges), and RL-based editing.

GraphRARE \cite{graphrare2024} is the most closely related prior work. It uses a deep
RL agent to select edges for addition or removal based on a \textit{node relative entropy}
metric that combines feature and structural information:
\begin{equation}
    E(v) = -\sum_{c} p(y=c|v) \log \frac{p(y=c|v)}{q(y=c|v)},
\end{equation}
where $p$ and $q$ are the predicted label distributions from different GNN layers.
The key differences between GraphRARE and GraphHARE are:

\begin{enumerate}
    \item \textbf{Reward signal.} GraphRARE's reward is derived from the relative entropy
    of label predictions; GraphHARE's reward explicitly combines task performance ($\Delta F_1$)
    and graph homophily ($\Delta h$), making the structural objective interpretable and tuneable.

    \item \textbf{GNN coupling.} GraphRARE jointly trains the RL agent and GNN backbone,
    requiring coordinated optimisation. GraphHARE freezes the backbone, separating graph
    editing from classification training and reducing computational cost.

    \item \textbf{Homophily awareness.} GraphRARE does not model homophily as an explicit
    objective. GraphHARE's $\beta$ parameter allows direct control over how much structural
    quality is traded against task performance.

    \item \textbf{Scalability.} Because GraphHARE operates on a frozen backbone, it can be
    applied post-hoc to any pre-trained GNN without retraining.
\end{enumerate}

These distinctions motivate the experimental comparison in Chapter~4, where GraphHARE is
evaluated alongside standard GNN, non-graph, autoencoder, and prior graph-editing baselines.
